\documentclass{article}
\usepackage{graphicx}
\usepackage{amsmath}
\usepackage{amssymb}
\usepackage{amsfonts}
\usepackage{hyperref}
\usepackage{exercises}
\usepackage{indentfirst}
\usepackage{circledsteps}
\usepackage{amsthm}
\usepackage{csquotes}
\usepackage{yhmath}
\usepackage{comment}
\usepackage{placeins}
\usepackage{geometry}
\usepackage{tikz}
\newcommand*{\oldneg}{\mathord{\sim}}
\newcommand{\Z}{\mathbb{Z}}
\newcommand{\N}{\mathbb{N}}
\let\oldemptyset\emptyset

\theoremstyle{definition}
\newtheorem{definition}{Definition}

\setlength{\parindent}{.5in}
\geometry{margin=1in}
\setlength{\parskip}{12pt}
\title{Chapter 1 Answers}
\begin{document}
\date{}
\maketitle

\section*{1.4}
Note that $P_1$ seems to be being used not just as an implication, but as an equivalence. See for example the proof of Lemma 4.3, where 
$\{a,b\}=\{a,d\}$ is said to imply $b\in \{a,d\}$. The equivalence is of course the definition that makes much more sense (if we only had the implication, we 
couldn't actually discern what $a=b$ even meant).

\subsection*{Induction with A4}
Consider the infinite sequence $a_0,a_1,a_2,...,a_n,...$ such that $a_0=\emptyset$ and for each $n$, $a_{n+1}=\{a_n\}$. We prove that 
if $n\ne m$, $a_n\ne a_m$. The base case is when $n=0$, that is, $a_n=\emptyset$. Obviously, $\emptyset\ne a_m$ for every $m\ne 0$, as every 
$a_m$ where $m>0$ is a set that contains an element, that is, a non-empty set. Now suppose the result holds for some $n$, that is, if $n\ne m$, then 
$a_n\ne a_m$. Now consider $a_{n+1}$. We prove that if $n+1\ne m$, then $a_{n+1}\ne a_m$. By similar reasoning as in the base case, $a_{n+1}\ne\emptyset$. 
Now consider any $m>0$. By Note 3, if $a_{n+1}=a_m$, then 
$a_n=a_{m-1}$, as $a_{n+1}=\{a_n\}$ and $a_m=\{a_{m-1}\}$. Now, if $n+1\ne m$, then $n\ne m-1$, and by the induction assumption, $a_n\ne a_{m-1}$, which is a contradiction. 
Thus $a_{n+1}\ne a_m$.
If $n+1=m$, the statement vacuously holds.

\subsection*{Theorem 4.4}
Note that this theorem can obviously be extended to equivalence, that is, $\langle a,b\rangle =\langle c,d\rangle$ iff $a=c$ and $b=d$.

\subsection*{$n$-tuples}
One can easily prove by induction that this scheme of $n$-tuples works. The base case is just Theorem 4.4. Now suppose the result holds for some $n$, that is, 
for all sets $x_1,...,x_n$ and $y_1,...,y_n$, if $\langle x_1,...,x_n\rangle =\langle y_1,...,y_n\rangle$ then for every $i\le n$, $x_i=y_i$. Now, if 
\[ \langle x_1,...,x_n,x_{n+1}\rangle =\langle y_1,...,y_n,y_{n+1}\rangle \] then 
\[ \langle \langle x_1,...,x_n\rangle ,x_{n+1}\rangle =\langle\langle y_1,...,y_n\rangle ,y_{n+1}\rangle \]
\[ \langle x_1,...,x_n\rangle =\langle y_1,...,y_n\rangle \]
and \[ x_{n+1}=y_{n+1} \] by Theorem 4.4 and then, by the induction assumption, for every $i\le n$,
$x_i=y_i$. Thus for all $i\le n+1$, $x_i=y_i$. Note that one could also take $\langle x_1,...,x_n,x_{n+1} \rangle$ to be 
$\langle x_1,\langle x_2,...,x_n,x_{n+1}\rangle\rangle$, as can be proven in a similar manner as above.


In fact, we can generalize the definitions of $n$-tuples in the following way. For $n=2$, define it in the same way as before. Then for every $n>2$, we could take 
$\langle x_1,...,x_n,x_{n+1}\rangle$ to be any $m$-tuple (where $m<n$) of the form 
$\langle x_1,...,x_a,\langle x_{a+1},...,x_b\rangle ,x_{b+1},...,x_c\rangle$ where, if $c>0$, then $c=n$, and if $c=0$, then $b=n$. 
This scheme can be proven to work by the following induction. Because 
$n=2$ is defined in the same way, the base case holds. Now suppose the result holds for all number less than $n$. Now consider any scheme for $n$ of the form:
\[ \langle x_1,...,x_a,\langle x_{a+1},...,x_b\rangle ,x_{b+1},...,x_c\rangle \] Suppose that 
\[ \langle x_1,...,x_a,\langle x_{a+1},...,x_b\rangle ,x_{b+1},...,x_c\rangle=\langle y_1,...,y_a,\langle y_{a+1},...,y_b\rangle ,y_{b+1},...,y_c\rangle \]
Because both $\langle x_1,...,x_a,\langle x_{a+1},...,x_b\rangle ,x_{b+1},...,x_c\rangle$ and $\langle y_1,...,y_a,\langle y_{a+1},...,y_b\rangle ,y_{b+1},...,y_c\rangle$ 
are some $m$-tuples where $m<n$, the induction assumption holds, and the three following things are true: (1) for all $i\le a$, $x_i=y_i$; 
(2) $\langle x_{a+1},...,x_b\rangle=\langle y_{a+1},...,y_b\rangle$; (3) for all $b+1\le i\le c$, $x_i=y_i$. By (2) and the induction assumption, 
($1'$) for all $a+1\le i\le b$, $x_i=y_i$. Thus, combining (1), ($1'$), and (3) together, we have, for all $i\le c$, $x_i=y_i$, that is, for all $i\le n$, 
$x_i=y_i$, thus completing the proof. 

\subsection*{Exercise 4.1}
(a): Suppose $\langle x,y\rangle$ is defined to be $\{\{\emptyset ,x\},\{\{\emptyset\},y\}\}$. Now suppose that 
$\langle a,b\rangle =\langle c,d\rangle$. Thus \[ \{\{\emptyset ,a\},\{\{\emptyset\},b\}\}=\{\{\emptyset ,c\},\{\{\emptyset\},d\}\} \]
meaning $\{\emptyset ,a\}\in \langle c,d\rangle$. Thus either $\{\emptyset ,a\}=\{\{\emptyset\},d\}$ or 
$\{\emptyset ,a\}=\{\emptyset ,c\}$. 


In the first case, we would have $\emptyset\in\{\{\emptyset\},d\}$, meaning $\emptyset=d$. Then also 
$a=\{\emptyset\}$. Then we also have two sub-cases: either $\{\{\emptyset\},b\}=\{\emptyset,\{c\}\}$ or 
$\{\{\emptyset\},b\}=\{\{\emptyset\},d\}$ (as $\{\{\emptyset\},b\}\in\langle c,d\rangle$). 
In the first sub-case, we would have $b=\emptyset$ and $c=\{\emptyset\}$, and because $a=\{\emptyset\}$ and 
$d=\emptyset$, the result would be proven. In the second sub-case we would have $\{\{\emptyset\},b\}=\{\{\emptyset\},d\}$ (meaning $b=d$ by Lemma 4.3) and thus 
\[ \{\{\emptyset ,a\},\{\{\emptyset\},b\}\}=\{\{\emptyset ,c\},\{\{\emptyset\},d\}\}=\{\{\emptyset ,c\},\{\{\emptyset\},b\}\} \]
that is,
\[ \{\{\{\emptyset\},b\},\{\emptyset ,a\}\}=\{\{\{\emptyset\},b\},\{\emptyset ,c\}\} \] and by Lemma 4.3 
$\{\emptyset ,a\}=\{\emptyset,c\}$, meaning $a=c$ (again by Lemma 4.3), thus proving the theorem.


In the second case, we have $\{\emptyset ,a\}=\{\emptyset ,c\}$ (meaning $a=c$ by Lemma 4.3) and thus $\{\{\emptyset\},b\}=\{\{\emptyset\},d\}$ by Lemma 4.3 
and, again by Lemma 4.3, $b=d$, thus completing the proof.


(b): Now suppose that $\langle x,y\rangle$ is defined to be $\{\{\emptyset,\{x\}\},\{\{y\}\}\}$. Now suppose that
$\langle a,b\rangle =\langle c,d\rangle$. Thus $\{\emptyset,\{a\}\}\in \langle c,d\rangle$, meaning either 
$\{\emptyset,\{a\}\}=\{\{d\}\}$ or $\{\emptyset,\{a\}\}=\{\emptyset,\{c\}\}$. The first sub-case is ruled out, as then we would have $\emptyset=\{\{d\}\}$, which is 
impossible, as $\emptyset$ contains no elements, whereas $\{\{d\}\}$ contains an element. Thus $\{\emptyset,\{a\}\}=\{\emptyset,\{c\}\}$ meaning, by Lemma 4.3, 
$\{a\}=\{c\}$ and thus $a=c$. Then again by Lemma 4.3, $\{\{b\}\}=\{\{d\}\}$ and thus $b=d$, hence showing $a=c$ and $b=d$.


Now consider the second case, where $\{\emptyset,\{a\}\}=\{\emptyset,\{c\}\}$. Thus $\{a\}=\{c\}$ (by Lemma 4.3) and $a=c$, meaning, again by Lemma 4.3, 
$\{\{b\}\}=\{\{d\}\}$ and thus $b=d$, therefore completing the proof.

\section*{1.5}

\subsection*{Exercise 5.1}
For suppose for any set $x$ there is a set $y$ that contains all elements of all elements of $x$. Now, $\cup x$ is defined to be the class of all elements of all elements 
of $A$. Thus $\cup x$ is a subclass of $y$, and hence, by $A_2$, is a set.

\subsection*{Exercise 5.2}
(a): Suppose $A\subseteq B$. Then, for all $x\in A$, $x\in B$. Now suppose $x\in \cup A$. Then $x\in y$ for some $y\in A$, that is, $y\in B$. Now, $\cup B$ is the set 
of all elements of all elements of $B$. But we have shown that $x$ is an element of $y$, an element of $B$. Thus $x\in \cup B$.
Hence $\cup A\subseteq B$.


(b): Suppose $A\subseteq B$. Then, for all $x\in A$, $x\in B$. Now suppose that $x\in\cap B$. Then $x\in y$ for every $y\in B$, and thus $x\in z$ for every 
$z\in A$. But $\cup A$ is the set of all elements of all elements of $A$, and we have shown that $x\in z$ for every element $z$ of $A$. Hence $x\in\cap A$. 
Thus $\cap B\subseteq\cap A$.

\subsection*{Exercise 5.3}
(a): This is true. For suppose that $A$ is transitive. Then all elements of all elements of $A$ are elements of $A$. $\cup A$ is the set of all elements of all elements of 
$A$. Now suppose that $x\in\cup A$. Thus, because it is an element of some element of $A$, we know that $x$ is an element of $A$ by the transitivity of $A$. Hence 
$x\in A$, and therefore $\cup A\subseteq A$.


(b): This is true. For suppose that $\cup A\subseteq A$. Then all elements of all elements of $A$ are elements of $A$, that is, $A$ is transitive.


(c): This is true. For suppose that $A$ is transitive. Then all elements of all elements of $A$ are elements of $A$. Now, $\cup A$ is the class of all elements of 
all elements of $A$. But suppose $\cup A$ was not transitive, that is, there was some element $x$ of some element $y$ of $\cup A$ such that $x\notin \cup A$. Now, we know that, 
because $y$ is an element of some element of $A$, and because $A$ is transitive, $y\in A$. But this means that $x$ is an element of an element of $A$, that is, 
$x\in\cup A$. Hence $\cup A$ is transitive. 


(d): This is true. For suppose that every element of $A$ is transitive. Then, for every element $x$ of $A$, all elements of all elements of $x$ are elements of $x$. 
Now suppose that $\cup A$ was not transitive. Hence there would be some element $x$ of an element $y$ of $\cup A$ such that $x\notin A$. Because $y\in\cup A$, 
$y$ is an element of some element $z$ of $A$. Because $z$ is transitive, all elements of all elements of $z$ are elements of $z$. But $x\in y$; hence $x\in z$. Thus 
$x$ is an element of an element of $A$, meaning $x\in\cup A$. Thus $\cup A$ is transitive. 

\subsection*{Exercise 5.4}
(a): On the one hand, suppose that $z\in x\cup y$. Then $z\in x$ or $z\in y$. To be an element of $\cup\{x,y\}$, $z$ must be an element of either $x$ or $y$. But we have 
shown this. Hence $z\in\cup\{x,y\}$. On the other hand, suppose that $z\in\cup\{x,y\}$. Then $z$ is either an element of $x$ or $y$. But this means $z\in x\cup y$. Hence, 
because $z\in x\cup y$ iff $z\in\cup\{x,y\}$, $x\cup y=\cup\{x,y\}$.


(b): On the ond hand, suppose that $z\in x\cap y$. Then $z\in x$ and $z\in y$. To be an element of $\cap\{x,y\}$, $z$ must be an element of both $x$ and $y$. But we have shown 
this. Hence $z\in\cap\{x,y\}$. On the other hand, suppose that $z\in\cap\{x,y\}$. Then $z\in x$ and $z\in y$, that is, $z\in x\cap y$. Hence 
$x\cap y=\cap\{ x,y\}$.

\subsection*{Exercise 5.5}
For non-empty sets, (a) and (b) are not necessarily true. In fact, assuming that sets cannot be members of themselves, they are never true. For suppose 
$\cup\{x\}=x$ and $x$ was non-empty. This would mean that $x$ was an element of $x$, which we assumed cannot occur. Similar reasoning shows that 
$\cap\{x\}\ne x$ for non-empty $x$. But if $x$ is empty, then both (a) and (b) are true. In particular, because $\emptyset$ has no elements, 
$\cup\{\emptyset\}=\cap\{\emptyset\}=\emptyset$.

\subsection*{Exercise 5.6}
(a): On the one hand, suppose that $x\in A\cap (B\cup C)$. Then $x\in A$ and $x\in B\cup C$. Hence either $x\in B$ or $x\in C$. If 
$x\in B$, then $x\in A\cap B$. If $x\in C$, then $x\in A\cap C$. Thus either $x\in A\cap B$ or $x\in A\cap C$, meaning
$x\in (A\cap B)\cup (A\cap C)$. On the other hand, suppose that $x\in (A\cap B)\cup (A\cap C)$. Then either $x\in A$ and $x\in B$, or 
$x\in A$ and $x\in C$. Thus $x\in A$ and either $x\in B$ or $x\in C$, that is, $x\in A\cap (B\cup C)$.


(b): On the one hand, suppose that $x\in A\cup (B\cap C)$. Then $x\in A$ or $x\in B\cap C$. If $x\in A$, then 
$x$ is in both $A\cup B$ and $A\cup C$, and hence also in $(A\cup B)\cap (A\cup C)$. If $x\in B\cap C$, then $x\in B$ and $x\in C$. 
Then $x\in A\cup B$ and $x\in A\cup C$, meaning $x\in (A\cup B)\cap (A\cup C)$. On the other hand, suppose that 
$x\in (A\cup B)\cap (A\cup C)$. Then $x\in A\cup B$ and $x\in A\cup C$. If $x\in A$, then $x\in A\cup (B\cap C)$. If $x\notin A$, then 
$x\in B$ and $x\in C$, that is, $x\in B\cap C$, meaning $x\in A\cup (B\cap C)$.


(c): On the one hand, suppose that $x\in A-(A-B)$. Then $x\in A$ and $x\notin A-B$. Thus either $x\notin A$ or $x\in B$. But it cannot be that $x\notin A$, for 
it was assumed that $x\in A$. Thus $x\in B$. Hence $x\in A$ and $x\in B$, meaning $x\in A\cap B$. On the other hand, suppose that $x\in A\cap B$. Then $x\in A$ and 
$x\in B$. Hence $x\notin A-B$. Thus $x\in A$ and $x\notin A-B$, meaning $x\in A-(A-B)$.


(d): The question gives the incorrect statement to prove. It should instead by $B-(A-B)=B$, not $\emptyset$. 
Now, $x\in B-(A-B)$ iff $x\in B$ and $x\notin A-B$. Thus either $x\notin A$ or $x\in B$. Thus this statement is satisfied for all $x\in B$. But 
suppose there were some $x\notin B$ such that $x\in B-(A-B)$. But that is impossible, as $x$ must be in $B$. Hence $x\in B$ iff $x\in B-(A-B)$, therefore meaning 
$B-(A-B)=B$.


(e): On the one hand, suppose that $A\subseteq B$. Now suppose that $x\in A\cup B$. Then $x\in A$. Now alternatively suppose that $x\in A$. 
But because $A\subseteq B$, $x\in B$. Thus $x\in A\cap B$. Thus $A\cap B=A$. On the other hand, suppose that $A\cap B=A$. Now suppose that $x\in A$. 
Thus $x\in A\cap B$, meaning $x\in B$. Hence $A\subseteq B$.


(f): On the one hand, suppose that $A\subseteq B$. Now suppose that $x\in A\cup B$. Then $x\in A$ or $x\in B$. If $x\in A$, then, because 
$A\subseteq B$, $x\in B$. In the second case, $x\in B$. Thus $x\in B$. Now alternatively suppose that $x\in B$. Thus $x\in A\cup B$. Thus 
$x\in A\cup B$ iff $x\in B$, meaning $A\cup B=B$. On the other hand, suppose that $A\cup B=B$. Now suppose that $x\in A$. Then 
$x\in A\cup B$, and thus $x\in B$.


(g): Suppose that $B\subseteq A$. By part (c) of this problem, $A-(A-B)=A\cap B$. By part (e) of this problem, $B\cap A=B$. But 
$B\cap A$ is just $A\cap B$. Thus $A\cap B=B$, meaning $A-(A-B)=B$.

\end{document}
